
\documentclass{article}
\usepackage{colortbl}
\usepackage{makecell}
\usepackage{multirow}
\usepackage{supertabular}

\begin{document}

\newcounter{utterance}

\centering \large Interaction Transcript for game `eqbench', experiment `eq\_bench', episode 158 with reflexion{-}Qwen3.5{-}9B.
\vspace{24pt}

{ \footnotesize  \setcounter{utterance}{1}
\setlength{\tabcolsep}{0pt}
\begin{supertabular}{c@{$\;$}|p{.15\linewidth}@{}p{.15\linewidth}p{.15\linewidth}p{.15\linewidth}p{.15\linewidth}p{.15\linewidth}}
    \# & \multicolumn{2}{c}{Player} && \multicolumn{2}{c}{Game Master} \\
    \hline

    \theutterance \stepcounter{utterance}  
    & & & \multicolumn{4}{p{0.6\linewidth}}{
        \cellcolor[rgb]{0.9,0.9,0.9}{
            \makecell[{{p{\linewidth}}}]{
                \texttt{\tiny{[P1$\langle$GM]}}
                \texttt{You have attempted this same task before and did not succeed. These are the plans you wrote after those attempts.} \\
\texttt{Trial \#1: New plan: I failed because I scored "Mocking" too high (8) and "Disappointed" too low (4), ignoring the context of their shared wager and Remy's specific admission of fear. Remy's tone shifted from sarcastic to serious when discussing the wager and his own vulnerability ("afraid to fail"). A score of 8 for Mocking ignores the gravity of the "fools together" moment, while a score of 4 for Disappointed underestimates the sting of Eilis calling him a "romantic fool" and his subsequent defensive pivot to "competition." Next time, I will prioritize "Disappointed" as the primary emotion (score 7{-}8) due to the rejection of his romantic overture and the framing of their relationship as purely competitive, while lowering "Mocking" to a secondary, residual tone (score 3{-}4) that reflects his initial sarcasm but not his current state. I will ensure the scores sum to a realistic emotional distribution where the most contextually supported emotion receives the highest value.} \\
\\ 
\texttt{Trial \#2: New plan: I failed because I misinterpreted Remy's shift from sarcasm to vulnerability. When Eilis called him a "romantic fool" and he pivoted to "competition," his primary emotion wasn't mocking; it was a defensive reaction to rejection, likely feeling **Disappointed** that his romantic overture ("here for you") was dismissed as charity. The "afraid to fail" line confirms a serious, internal struggle, not a lighthearted mockery. I should have scored **Disappointed** highest (8{-}9) to reflect the sting of rejection and the defensive hardening of his attitude, while keeping **Mocking** low (2{-}3) as a residual trait from his earlier banter, and **Empathetic** at zero since he explicitly rejected her need for charity. I will ensure the scores reflect this dynamic: Disappointment as the dominant response to the failed romantic connection, with Mocking as a minor, fading element.} \\
\\ 
\texttt{Trial \#3: New plan: I failed because my previous attempts incorrectly prioritized "Mocking" and "Astonished" while underestimating the depth of "Disappointed." Remy's shift from playful banter to a serious wager, followed by his admission of being "afraid to fail" after Eilis dismissed his romantic overture ("here for you"), indicates a profound sense of rejection and unrequited effort. The phrase "romantic fool" stings, and his pivot to "competition" is a defensive mechanism masking this disappointment. I must re{-}evaluate the scores to reflect that "Disappointed" is the dominant emotion (9{-}10) due to the failed connection and the harsh reality of the wager, while "Mocking" should be low (1{-}2) as it is merely a fading trait, and "Astonished" and "Empathetic" should be zero or negligible given the adversarial and defensive tone of the final exchange.} \\
\\ 
\\ 
\\ 
\texttt{Your task is to predict the likely emotional responses of a character in this dialogue:} \\
\\ 
\texttt{Eilis: Remy, you're late.} \\
\texttt{Remy: And you're as punctual as ever, Eilis. Still chasing shadows, are we?} \\
\texttt{Eilis: Better shadows than illusions of grandeur, Remy.} \\
\texttt{Remy: Ah, the sting of your wit. I've missed that. } \\
\texttt{Eilis: And I've missed your incessant need to prove something.} \\
\texttt{Remy: Is that so? Well, I suppose we're both here for the same reason then.} \\
\texttt{Eilis: Don't flatter yourself, Remy. Our reasons couldn't be more different.} \\
\texttt{Remy: Oh, do enlighten me, Eilis. } \\
\texttt{Eilis: You're here to feed your ego. I'm here to reclaim what I lost.} \\
\texttt{Remy: And what if I told you I'm here for you, Eilis?} \\
\texttt{Eilis: I'd say you're a bigger fool than I thought.} \\
\texttt{Remy: Well, then let's be fools together, shall we? } \\
\texttt{Eilis: I don't need your charity, Remy.} \\
\texttt{Remy: It's not charity, Eilis. It's competition. } \\
\texttt{Eilis: Is that a challenge, Remy?} \\
\texttt{Remy: If you're not afraid, Eilis.} \\
\texttt{Eilis: Afraid? Of you? Don't make me laugh.} \\
\texttt{Remy: Then let's make a wager, Eilis. } \\
\texttt{Eilis: A wager? And what would that be?} \\
\texttt{Remy: The first to create the Philosopher's Stone.} \\
\texttt{Eilis: And what does the loser get, Remy?} \\
\texttt{Remy: The satisfaction of knowing they gave it their all, Eilis.} \\
\texttt{Eilis: You always were a romantic fool, Remy.} \\
\texttt{Remy: And you were always afraid to lose, Eilis.} \\
\texttt{Eilis: I'm not afraid to lose, Remy. I'm afraid to fail.} \\
\texttt{{[}End dialogue{]}} \\
\\ 
\texttt{At the end of this dialogue, Remy would feel...} \\
\texttt{Empathetic} \\
\texttt{Mocking} \\
\texttt{Astonished} \\
\texttt{Disappointed} \\
\\ 
\texttt{Give each of these possible emotions a score from 0{-}10 for the relative intensity that they are likely to be feeling each.} \\
\\ 
\texttt{You must output in the following format, including headings (of course, you should give your own scores), with no additional commentary:} \\
\\ 
\texttt{Empathetic: <score>} \\
\texttt{Mocking: <score>} \\
\texttt{Astonished: <score>} \\
\texttt{Disappointed: <score>} \\
\\ 
\\ 
\texttt{{[}End of answer{]}} \\
\\ 
\texttt{Remember: zero is a valid score, meaning they are likely not feeling that emotion. You must score at least one emotion > 0.} \\
\\ 
\texttt{Your answer:} \\
            }
        }
    }
     \\ \\

    \theutterance \stepcounter{utterance}  
    & \multicolumn{4}{p{0.6\linewidth}}{
        \cellcolor[rgb]{0.561,0.737,0.561}{
            \makecell[{{p{\linewidth}}}]{
                \texttt{\tiny{[P1$\rangle$GM]}}
                \texttt{Empathetic: 0} \\
\texttt{Mocking: 2} \\
\texttt{Astonished: 0} \\
\texttt{Disappointed: 9} \\
            }
        }
    }
    & & \\ \\

    \theutterance \stepcounter{utterance}  
    & & & \multicolumn{2}{p{0.3\linewidth}}{
        \cellcolor[rgb]{0.9,0.9,0.9}{
            \makecell[{{p{\linewidth}}}]{
                \texttt{\tiny{[GM$|$GM]}}
                \texttt{Empathetic: 0} \\
\texttt{Mocking: 2} \\
\texttt{Astonished: 0} \\
\texttt{Disappointed: 9} \\
            }
        }
    }
    & & \\ \\

    \theutterance \stepcounter{utterance}  
    & & & \multicolumn{2}{p{0.3\linewidth}}{
        \cellcolor[rgb]{0.9,0.9,0.9}{
            \makecell[{{p{\linewidth}}}]{
                \texttt{\tiny{[GM$|$GM]}}
                \texttt{Empathetic: 6} \\
\texttt{Mocking: 7} \\
\texttt{Astonished: 0} \\
\texttt{Disappointed: 2} \\
            }
        }
    }
    & & \\ \\

    \theutterance \stepcounter{utterance}  
    & & & \multicolumn{2}{p{0.3\linewidth}}{
        \cellcolor[rgb]{0.9,0.9,0.9}{
            \makecell[{{p{\linewidth}}}]{
                \texttt{\tiny{[GM$|$GM]}}
                \texttt{game\_result = LOSE} \\
            }
        }
    }
    & & \\ \\

    \theutterance \stepcounter{utterance}  
    & & \multicolumn{4}{p{0.6\linewidth}}{
        \cellcolor[rgb]{0.95,0.95,0.95}{
            \makecell[{{p{\linewidth}}}]{
                \texttt{\tiny{[P1$\leftrightarrow$P]}}
                \texttt{Post{-}episode Reflexion (not sent to GM):} \\
\texttt{New plan: I failed because I incorrectly identified "Mocking" as the dominant emotion, ignoring the narrative shift from sarcasm to serious vulnerability. Remy's admission of being "afraid to fail" and his defensive pivot to "competition" after being called a "romantic fool" indicate that his primary feeling is **Disappointment** regarding the rejection of his romantic overture. The wager is a mask for this hurt, not an expression of mockery. I must re{-}score to reflect **Disappointed** as the highest intensity (9{-}10), representing the sting of rejection and the defensive hardening of his attitude. **Mocking** should be reduced to a low residual score (1{-}2) reflecting only his initial tone, while **Empathetic** and **Astonished** remain at zero, as the dialogue is adversarial and grounded in personal failure rather than surprise or shared feeling.} \\
            }
        }
    }
    & \\ \\

\end{supertabular}
}

\end{document}
